\documentclass[../../../include/open-logic-section]{subfiles}

\begin{document}

\olfileid{sth}{story}{predicative}

\olsection{स्वसमावेशरहित आणि स्वसमावेशक व्याख्या}

रसेलचा संच~$R$ हा
$\Setabs{x}{x \notin
x}$ वापरून परिभाषित
केला होता. अधिक
स्पष्टपणे सांगायचे
तर~$R$ म्हणजे
स्वतःचे सदस्य
नसलेले, आणि
फक्त असेच,
सर्व संच सदस्य
म्हणून असलेला
संच ठरला असता.
म्हणून~$R$ ची
व्याख्या करताना
आपण~$R$
(अस्तित्वात असता)
ज्या प्रांतात
आला असता
त्या प्रांतावर
परिमाणन करतो.

ही \emph{स्वसमावेशक}
व्याख्या आहे.
सर्वसाधारणपणे,
परिभाषित होणारी
वस्तू ज्या
प्रांतात समाविष्ट
आहे त्या
प्रांतावर
परिमाणन करणारी
व्याख्या तेव्हा
आणि तेव्हाच
स्वसमावेशक
असते असे
म्हणता येईल.

या विरोधापत्तींच्या
पार्श्वभूमीवर
व्हाइटहेड, रसेल,
प्वाँकारे आणि
वाइल यांनी
अशा स्वसमावेशक
व्याख्यांना
‘‘दुष्टचक्र निर्माण
करणाऱ्या’’ म्हणून
नाकारले:
\begin{quote}
  टाळायच्या
  विरोधापत्तींचे
  विश्लेषण
  दाखवते की
  त्या सर्व
  एका प्रकारच्या
  दुष्टचक्रातून
  निर्माण होतात.
  एखाद्या
  वस्तुसंग्रहात
  असे सदस्य
  असू शकतात
  जे फक्त
  त्या संपूर्ण
  संग्रहाच्या
  साहाय्यानेच
  परिभाषित
  करता येतात,
  असे मानल्याने
  संबंधित दुष्टचक्रे
  तयार होतात
  [\ldots. \textparagraph]

  अनुचित समग्रता
  टाळू देणारे
  तत्त्व असे
  मांडता येते:
  ‘संग्रहातील
  \emph{सर्वांचा}
  उल्लेख ज्यात
  येतो, ती
  वस्तू त्या
  संग्रहातील
  एक वस्तू
  असू नये’;
  किंवा उलट:
  ‘एखाद्या
  संग्रहाला
  समग्रता असती
  तर त्यातील
  काही सदस्य
  फक्त त्या
  समग्रतेच्या
  साहाय्यानेच
  परिभाषित
  करता आले
  असते, तर
  त्या संग्रहाची
  समग्रता
  अस्तित्वात
  नाही.’ याला
  आपण
  ‘दुष्टचक्र
  तत्त्व’
  म्हणू,
  कारण
  अनुचित
  समग्रता
  गृहीत
  धरल्याने
  उद्भवणारी
  दुष्टचक्रे
  ते टाळू
  देते.
  \citep[p.~37]{WhiteheadRussell1910}
\end{quote}
त्यांच्याप्रमाणे
\emph{दुष्टचक्र तत्त्व}
स्वीकारले, तर
\olref[sth][story][rus]{sec}
मधील विनाशकारी
अनौपचारिक
गुणधर्माधारित
संचरचना-रूपबंधाऐवजी
पुढील प्रकारचे
तत्त्व वापरण्याचा
प्रयत्न करता
येईल:

\begin{defish}
\emph{स्वसमावेशरहित
गुणधर्माधारित
संचरचना.}
फक्त संचांवर
परिमाणन करणाऱ्या
प्रत्येक सूत्र~$\phi$
साठी:
संच$^\prime$
$\Setabs{x}{\phi(x)}$
अस्तित्वात असतो.
\end{defish}

संच$^{\prime}$
हे साधे संच
नसतील, तोवर
व्याघात उद्भवणार
नाही.

दुर्दैवाने,
स्वसमावेशरहित
गुणधर्माधारित
संचरचना फारशी
\emph{व्यापक}
नाही. कारण
तिच्यामुळे
संच$^\prime$
या नव्या
वस्तू मिळतात.
म्हणून संच$^\prime$
यांच्यावर
परिमाणन करणारी
सूत्रेही विचारात
घ्यावी लागतील.
ती नेहमी
संच$^\prime$
देत असतील,
तर स्वतःचे
सदस्य नसलेल्या
सर्व संच$^\prime$
यांचा संच$^\prime$
पाहताच रसेलची
विरोधापत्ती
पुन्हा उद्भवेल.
त्यामुळे याच
विचारानुसार
संच$^\prime$
यांच्यावर
परिमाणन करणारे
सूत्र अनुरूप
संच$^{\prime\prime}$
देते असे
म्हणावे लागेल.
पुढे
संच$^{\prime\prime\prime}$,
संच$^{\prime\prime\prime\prime}$
इत्यादी लागतील.
अशी कितीतरी
प्राइम चिन्हे
टाळण्यासाठी
यांना संच$_0$,
संच$_1$,
संच$_2$,
संच$_3$,
संच$_4$,\ldots
असे मानणे
सोपे ठरेल.
यातून (सरळ)
प्रकार-उपपत्तीचा
मार्ग मिळतो.

अशा उपपत्तीवर
काही उघड
आक्षेप आहेत
(मात्र ते
\emph{निर्णायक}
आहेतच असे
स्पष्ट नाही).
थोडक्यात:
तिचा वापर
जड जातो;
ती अनेक
भिन्न प्रकारच्या
वस्तू गृहीत
धरते; आणि
शेवटी
स्वसमावेशक
व्याख्या
\emph{खरोखर
इतक्या वाईट}
आहेत का,
हेही स्पष्ट
नाही.

शेवटचा मुद्दा
स्पष्ट करण्यासाठी
\citeauthor{Ramsey1925}
यांची पुढील
टिप्पणी पाहा:
\begin{quote}
  एखाद्या
  गटातील
  सर्वांत
  उंच माणूस
  म्हणून आपण
  एखाद्याचा
  उल्लेख करू
  शकतो. असे
  केल्याने तो
  स्वतः ज्या
  समग्रतेचा
  सदस्य आहे
  तिच्याच
  साहाय्याने
  त्याची ओळख
  पटते; तरीही
  त्यात दुष्टचक्र
  नसते.
  \citep{Ramsey1925}
\end{quote}
रॅम्झी यांचा
मुद्दा असा
की ‘‘गटातील
सर्वांत उंच
माणूस’’ ही
\emph{स्वसमावेशक}
व्याख्या आहे;
पण ती
स्पष्टपणे
निर्दोष आहे.

त्यावर असे
उत्तर देता
येईल की
या उदाहरणात
सर्वांत उंच
व्यक्तीला
\emph{स्वसमावेशरहित}
मार्गानेही
ओळखता येईल.
उदाहरणार्थ,
आपण त्या
माणसाकडे
सरळ बोट
दाखवू
शकतो. मग
स्वसमावेशक
व्याख्यांवरील
आक्षेप फक्त
\emph{स्वसमावेशक
मार्गानेच}
ओळखता येणाऱ्या
वस्तूंपुरता
मर्यादित करावा
लागेल. पण
अशा
‘‘मूलभूत
स्वसमावेशकतेत’’
नेमके वाईट
काय आहे,
हेही पुढे
स्पष्ट करावे
लागेल.\footnote{
  अधिक माहितीसाठी
  \citet{Linnebo2010}
  पाहा.}

वस्तू
\emph{निर्माण}
करण्यासाठी
आपल्या
व्याख्यांनी
कृतीसूचना
द्यावी असे
अपेक्षित
असेल, तर
स्वसमावेशक
व्याख्या
निश्चितच
मोठी अडचण
ठरतात.
कारण अशा
व्याख्येवरून
वस्तू बनवताना
खरेच
दुष्टचक्रात
अडकावे
लागेल:
स्वसमावेशकपणे
निर्दिष्ट
वस्तू निर्माण
करण्यासाठी
\emph{आधी} सर्व
वस्तू
(त्यात ती
निर्दिष्ट
वस्तूही येते)
निर्माण कराव्या
लागतील,
कारण तिची
व्याख्या
सर्व वस्तूंच्या
संदर्भात आहे.
म्हणजे त्या
वस्तूची
निर्मिती
होण्यापूर्वीच
तिची निर्मिती
करावी लागेल.
पण पुन्हा,
‘‘मूलभूत
स्वसमावेशक’’
रीतीने
निर्दिष्ट
संचांवरील
हा गंभीर
आक्षेप
तेव्हाच
ठरतो,
जेव्हा
संच म्हणजे
आपण
\emph{निर्माण}
करत असलेल्या
वस्तू आहेत
असे मानतो.
आपण
(बहुधा)
तसे मानत
नाही.

त्यामुळे,
चांगल्या किंवा
वाईट कारणांनी,
रूढ झालेल्या
पद्धतीत
(अ)स्वसमावेशकतेबाबत
कठोर भूमिका
घेतलेली नाही.
त्याऐवजी
आता संचयी-क्रमिक
पद्धत मानली
जाते. शेवटी
तिच्यामुळे
संचांचे
‘‘टप्प्यां’’मध्ये
स्तरीकरण करता
येते---स्वसमावेशरहित
पद्धतीत वस्तू
संच$_0$,
संच$_1$,
संच$_2$,
\ldots मध्ये
स्तरीकृत
होतात तसेच
\emph{थोड्याफार}
प्रमाणात.
मात्र या
टप्प्यांतील
संचांच्या
मूलभूत प्रकारांत
कोणताही
फरक गृहीत
धरला जात
नाही.

\end{document}
